\begin{vidu}
Mạch chọn sóng của một máy thu sóng vô tuyến gồm cuộn cảm thuần có độ tự cảm $\dfrac{0,4}{\pi}$ H và tụ điện có điện dung $C = \dfrac{10}{4\pi}$ pF.
\begin{itemize}
\item[a)] Tính tần số góc, chu kì dao động riêng, tần số dao động riêng của mạch.
\item[b)] Giữ nguyên độ tự cảm L và thay tụ C bằng tụ điện có điện dung biến đổi từ 10 pF đến 640 pF. Khi đó chu kì dao động riêng của mạch này có giá trị thay đổi như thế nào? 
\item[c)] Điện dung tụ điện ở câu b) được thay đổi bằng cách xoay một bản tụ để thay đổi phần diện tích đối diện hai bản tụ. Hỏi phần diện tích này phải thay đổi như thế nào?
\item[d)] Thay tụ điện C bằng tụ C' thì thấy chu kì dao động riêng của mạch bằng $4.10^{-6}\ s$. Tính chu kì dao động riêng của mạch khi ghép hai tụ điện C và C' nối tiếp, song song.
\end{itemize}
\loigiai{
\begin{itemize}
\item[a)] 	Chu kì dao động riêng của mạch:
\begin{align*}
T=2\pi \sqrt{LC}=2\pi \sqrt{\dfrac{0,4}{\pi}\dfrac{10}{4\pi}.10^{-12}}=2.10^{-6}\ s.
\end{align*}\\
Tần số góc của mạch: $\omega=\dfrac{2\pi}{T}=\pi.10^{6}\ rad/s$.\\
Tần số dao động riêng của mạch: $f=\dfrac{1}{T}=5.10^5\ Hz$.			
\item[b)] Theo công thức tính chu kì dao động của mạch: $T=2\pi\sqrt{LC}
\Rightarrow 2\pi\sqrt{L{C_{\min}}}\le T\le 2\pi\sqrt{L{C_{\max}}}$.\\
Thay số vào ta được: $7,09.10^{-6}\ s \leq T \leq 5,76.10^{-5}\ s$.
\item[c)] Ta có: $C=\dfrac{\epsilon S}{4\pi kd}\Rightarrow C\sim S \Rightarrow$ C tăng 64 lần nên S tăng 64 lần.
\item[d)] Do $T^2\sim C 
\begin{cases}
C_{ss}=C_1+C_2 \Rightarrow T_{ss}^2=T_1^2+T_2^2\Rightarrow T_{ss}=\sqrt{T_1^2+T_2^2}=4,47.10^{-6}\ s.\\
\dfrac{1}{C_{nt}}=\dfrac{1}{C_1}+\dfrac{1}{C_2}\Rightarrow \dfrac{1}{T_{nt}^2}=\dfrac{1}{T_1^2}+\dfrac{1}{T_2^2}\Rightarrow T_{nt}\approx 1,79.10^{-6}\ s.
\end{cases}$ 
\end{itemize}
}
\end{vidu}
\loai{Tính tần số góc}
\begin{vidu}%[3P4-1.2B][Loai1]
Mạch dao động LC gồm tụ $C = 16$ nF và cuộn cảm $L = 25$ mH. Tần số góc dao động của mạch là
\choice
{$\omega  = 2000$ rad/s}
{$\omega  = 200$ rad/s}
{\True $\omega  = 5.10^4$ rad/s}
{$\omega  = 5.10^{-4}$ rad/s}
\loigiai{
Ta có: $\omega =\dfrac{1}{\sqrt{LC}}=\dfrac{1}{\sqrt{16.10^{-9}.25.10^{-3}}}=5.10^4\ rad/s$. 
}
\end{vidu}
\loai{Tính chu kì}
\begin{vidu}%[3P4-1.2B][Loai2]
Mạch dao động LC lí tưởng có tụ điện $C = 25$ pF và cuộn cảm $L=4.10^{-4}\ H$. Chu kì dao động của mạch là
\choice
{\True $2\pi .10^{-7}$ s}
{$10^7$ s}
{$2.10^{-7}$ s}
{$10^7$  rad/s}
\loigiai{
Chu kì của mạch LC: $T=2\pi \sqrt{LC}=2\pi.10^{-7}$ s.
}
\end{vidu}
\loai{Tính tần số}
\begin{vidu}%[3P4-1.2B][Loai3] 
Mạch dao động LC có điện tích trong mạch biến thiên điều hòa theo phương trình $q=\cos({2\pi .10^4t})\ \mu C$. Tần số dao động của mạch là
\choice
{$f=10\ Hz$}
{\True $f=10\ kHz$}
{$f=2\pi \ Hz$}
{$f=2\pi \ Hz$}
\loigiai{
Ta có: $f=\dfrac{\omega}{2\pi}=10000\ Hz=10\ kHz$.
}
\end{vidu}
\loai{Tính L, C}
\begin{vidu}%[3P4-1.2B][Loai4]
Một mạch dao động LC gồm một cuộn cảm có độ tự cảm $L = \dfrac{1}{2\pi}$ H và một tụ điện có điện dung C. Tần số dao động riêng của mạch là 1 MHz. Giá trị của C bằng
\choice
{\True $C = \dfrac{1}{2\pi}$ pF}
{$C = \dfrac{1}{2\pi}$ F}
{$C = \dfrac{1}{8\pi}$ mF}
{$C = \dfrac{1}{6\pi}\ \mu F$}
\loigiai{
Ta có: $f=10^{6}=\dfrac{1}{2\pi \sqrt{LC}}\Rightarrow LC=\dfrac{1}{10^{12}.4{\pi}^2}\Rightarrow C=\dfrac{1}{2\pi}.10^{-12}\ F=\dfrac{1}{4\pi}\ pF$. 
}
\end{vidu}
\loai{Thay đổi L, C}
\begin{vidu}%[3P4-1.2K][Loai5] 
\nguon{ĐH - 2010} Một mạch dao động lí tưởng gồm cuộn cảm thuần có độ tự cảm 4 $\mu H$ và một tụ điện có điện dung biến đổi từ 10 pF đến 640 pF. Lấy $\pi ^2 = 10$. Chu kì dao động riêng của mạch này có giá trị
\choice
{từ $2.10^{-8}$ s đến $3,6.10^{-7}$ s}
{từ $4.10^{-8}$ s đến $2,4.10^{-7}$ s}
{\True từ $4.10^{-8}$ s đến $3,2.10^{-7}$ s}
{từ $2.10^{-8}$ s đến $3.10^{-7}$ s}
\loigiai{
\begin{itemize}
\item Theo công thức tính chu kì dao động của mạch: $T=2\pi\sqrt{LC}
\Rightarrow 2\pi\sqrt{L{C_{\min}}}\le T\le 2\pi\sqrt{L{C_{\max}}}$.
\item Thay số vào ta được $4.10^{-8}$ s đến $3,2.10^{-7}$ s.
\end{itemize}
}
\end{vidu}
\loai{Cấu tạo tụ điện}
\begin{vidu}%[3P4-1.2K][Loai6]
Một mạch dao động LC lí tưởng có thể biến đổi trong dải tần số từ 10 MHz đến 50 MHz bằng cách thay đổi khoảng cách giữa hai bản tụ điện phẳng. Khoảng cách giữa các bản tụ thay đổi
\choice
{5 lần}
{16 lần}
{160 lần}
{\True 25 lần}
\loigiai{
Ta có: $\dfrac{f_2}{f_1}=\dfrac{\dfrac{1}{2\pi \sqrt{LC_2}}}{\dfrac{1}{2\pi \sqrt{LC_1}}}=\sqrt{\dfrac{C_1}{C_2}}=\sqrt{\dfrac{d_2}{d_1}}\Rightarrow \dfrac{d_2}{d_1}=\left(\dfrac{f_2}{f_1}\right)^2=25$.
}
\end{vidu}
\loai{Ghép tụ}
\begin{vidu}%[3P4-1.2K][Loai7] 
Mạch dao động lí tưởng gồm cuộn dây thuần cảm L không đổi, tụ điện có điện dung C thay đổi. Khi $C=C_1$ thì mạch dao động với tần số $30\ MHz,$ khi $C=C_1+C_2$ thì mạch dao động với tần số $24\ MHz$. Khi $C=4C_2$ thì mạch dao động với tần số là
\choice
{\True $20\ MHz$}
{$80\ MHz$}
{$40\ MHz$}
{$60\ MHz$}
\loigiai{
\begin{itemize}
\item Tần số của mạch dao động LC được tính theo công thức $f=\dfrac{1}{2\pi \sqrt{LC}}\Rightarrow f\sim \dfrac{1}{\sqrt{C}}\Leftrightarrow f^2\sim \dfrac{1}{C}$.
\item Khi $C=C_1$ thì $f_1^2\sim \dfrac{1}{C_1}\Rightarrow C_1\sim \dfrac{1}{f_1^2}$.
\item Khi $C=C_1+C_2$ thì $f_{12}^2\sim \dfrac{1}{C_1+C_2}\Rightarrow C_1+C_2\sim \dfrac{1}{f_{12}^2}$.
\item Từ đó suy ra $4C_2\sim 4\left(\dfrac{1}{f_{12}^2}-\dfrac{1}{f_1^2}\right)\Leftrightarrow \dfrac{1}{f^2}=4\left(\dfrac{1}{f_{12}^2}-\dfrac{1}{f_1^2}\right)$.
\item Thay số vào ta tính được tần số khi $C=4C_2$ là $f=20\ MHz$.
\end{itemize}}
\end{vidu}
\begin{vidu}%[3P4-1.2K][Loai7]
Một mạch dao động điện từ có cuộn cảm không đổi L. Nếu thay tụ điện C bởi các tụ điện $C_1,\ C_2,\ C_1$ nối tiếp $C_2$ và $C_1$ song song $C_2$ thì chu kì dao động riêng của mạch lần lượt là $T_1,\ T_2,\ T_{nt} = 4,8\ \mu s,\ T_{ss} = 10\ \mu s$. Hãy xác định $T_1$, biết $T_1 > T_2$?
\choice
{$T_1 = 9\ \mu s$}
{\True $T_1 = 8\ \mu s$}
{$T_1 = 10\ \mu s$}
{$T_1 = 6\ \mu s$}
\loigiai{
\begin{itemize}
\item Ta có $T=2\pi \sqrt{LC}\Rightarrow T^2\sim C$.
\item Khi hai tụ điện mắc song song: $T^2=T_1^2+T_2^2\Leftrightarrow T_1^2+T_2^2=10^2\quad (1)$.
\item Khi hai tụ điện mắc nối tiếp: $\dfrac{1}{T^2}=\dfrac{1}{T_1^2}+\dfrac{1}{T_2^2}\Leftrightarrow \dfrac{1}{T_1^2}+\dfrac{1}{T_2^2}=\dfrac{1}{4,8^2}\quad (2)$.
\item Từ $(1)$ và $(2)$ suy ra $T_1=8\ \mu s,\ T_2=6\ \mu s$.
\end{itemize}}
\end{vidu}
\loai{Nạp năng lượng}
\begin{vidu}%[3P4-1.2K][Loai8] 
Đoạn mạch AB gồm cuộn cảm thuần nối tiếp với tụ điện. Đặt nguồn xoay chiều có tần số góc $\omega $ vào hai đầu A và B thì tụ điện có dung kháng $100\ \Omega $, cuộn cảm có cảm kháng $50\ \Omega $. Ngắt A, B ra khỏi nguồn và tăng độ tự cảm của cuộn cảm một lượng 0,5 H rồi nối A và B thành mạch kín thì tần số góc dao động riêng của mạch là 100  rad/s. Tính $\omega $.
\choice
{$80\pi \ rad/s$}
{$50\pi \ rad/s$}
{\True $100\ rad/s$}
{$50\ rad/s$}
\loigiai{
Ta có: 
$\begin{cases}
Z_L=\omega L=50\ \Omega \Rightarrow L=\dfrac{50}{\omega}\\
Z_C=\dfrac{1}{\omega C}=100\ \Omega \Rightarrow C=\dfrac{1}{100\omega}\\
\dfrac{1}{\omega_0^2}=L'C=\left(L+\Delta L\right)C.
\end{cases}
\Rightarrow \dfrac{1}{10000}=\dfrac{50}{\omega}\dfrac{1}{100\omega}+0,5.\dfrac{1}{100\omega}\Rightarrow \omega =100 \ rad/s$.
}
\end{vidu}
